\section{Conceptual framework}

\subsection{Engagement as a situated process}

Kahu \citep{kahu2013} describes student engagement as a complex, multifaceted construct that cannot be understood adequately from a single behavioural or psychological perspective. The framework separates sociocultural context, structural and psychosocial influences, engagement itself, and proximal and distal consequences. This distinction matters empirically because many variables routinely described as engagement are better understood as conditions that shape engagement or as outcomes that follow from it. For administrative research, the risk is particularly acute: an analyst can observe that a student visited a support centre, failed a course, changed instructor, or progressed to a later course, but those observations do not reveal the student's affective or cognitive state.

The present study therefore uses Kahu's framework as a \emph{theoretical organizer}, not as a claim that administrative records measure latent engagement. We treat the assigned MU classroom and degree programme as structural context, classroom-relative performance and course outcomes as academic experience and feedback, CMAT attendance as a formal help-seeking trace, later instructor enrolment as a context-selection trace, and subsequent performance or progression as outcomes. The value of this separation is that different observed behaviours can move in different directions without forcing them onto a single engagement scale.

This perspective is also consistent with broader research on the first-year experience. Students' adjustment to higher education involves not only formal academic performance but also new relationships, sources of support, and institutional environments. Qualitative work on first-year students has shown that academic and social integration emerge through multiple interacting processes rather than a single behavioural index \citep{wilcox2005}. Our data do not observe the social dimensions examined in that literature, which is precisely why we avoid treating CMAT use as an exhaustive measure of engagement. Instead, formal support use is one institutional trace within a broader system of possible responses.

\subsection{Help-seeking and the problem of non-engagement}

Academic help-seeking is commonly framed as a self-regulated response to difficulty, but the behaviour is heterogeneous. Fong et al.'s meta-analysis \citep{fong2023} distinguishes formal from informal sources and instrumental from executive or avoidant forms, showing that the association with achievement depends on what kind of help is sought and for what purpose. Formal support-centre attendance is therefore meaningful but incomplete: it indicates contact with an institutional support resource, not the quality of the interaction, the student's motivation, or the availability of alternative forms of help.

Mathematics support research has long confronted the same selection problem. Support centres provide additional assistance outside regular teaching \citep{lawson2020,mullen2024}, and observational studies often find more favourable outcomes among attendees after some statistical adjustment \citep{macanbhaird2009,jacob2019}. At the same time, students who attend differ from those who do not, and evaluation is complicated by unobserved study effort, prior preparation, perceived need, and voluntary uptake. Mac an Bhaird et al. \citep{macanbhaird2013} document that a meaningful group of students do not engage appropriately with available mathematics learning supports, despite the apparent benefits of such provision. This non-engagement problem makes it unsafe to assume that the students in greatest academic difficulty are necessarily the students most likely to seek formal assistance.

Recent syntheses reinforce the need for caution. Mullen et al. \citep{mullen2024} identify a broad and predominantly positive mathematics and statistics support literature, but also emphasize methodological diversity, potential bias, and the need for more rigorous evaluation designs. The present paper does not attempt to identify a causal impact of CMAT attendance. Instead, it focuses on the pattern of formal help-seeking itself and on how that pattern coexists with academic challenge and later adaptation.

\subsection{Academic challenge, feedback, and adaptation}

A longitudinal perspective allows academic difficulty to be treated as feedback that may precede later behavioural change, while preserving the distinction between response and outcome. Students who experience difficulty may seek formal help, intensify independent study, rely on peers, change schedules, select different instructors, repeat a course, or disengage from the academic pathway. Administrative data observe only some of these responses, but repeated observations can still reveal whether support use and instructional context change after salient academic events.

Two transitions are especially informative here. The first is the move from MU to Calculus I. Because MU sections are institutionally assigned while later Calculus enrolment occurs within a set of offered sections, the transition permits a limited analysis of context selection. The second is the transition from a failed MU attempt to the next MU attempt. Failure is an unambiguous academic event that precedes the next enrolment, allowing us to ask whether students subsequently change instructor context or support use.

Neither transition identifies intention. A student who enrols with an instructor associated with historically higher pass rates may have actively sought that instructor, may have selected a convenient schedule, or may have been constrained by capacity or registration rules. Accordingly, throughout the paper we describe \emph{subsequent enrolment with historically higher- or lower-outcome instructors}, rather than deliberate choice of ``easy'' or ``hard'' professors.

